\section[Critères structurels et carte métrologique de \(G\)]{Critères de
sélection de la coordonnée gravitationnelle : holonomie enrichie, réponse
spectrale et carte métrologique}
\label{sec:selector-gravitatorio-rev11}
\index{constante gravitationnelle!transducteur et sélecteur}
\index{gravité!coordonnée métrologique}
\index{holonomie!sélecteur gravitationnel}

La composition APP--TRIT--TPK, l'état enrichi et la branche d'action
construisent d'abord la coordonnée gravitationnelle. La constante gravitationnelle
s'organise en trois niveaux typés : la droite dimensionnelle qui abrite \(G\), le
transducteur qui convertit une coordonnée adimensionnelle en une section de cette
droite et le sélecteur structurel de cette coordonnée. L'action
Palatini--Cartan--Holst et son hamiltonien contraint réalisent ensuite la
sortie déjà générée. La carte décimale et l'incertitude expérimentale sont
présentées à la fin comme représentation et confrontation.

\subsection{Droite dimensionnelle et transducteur interne}

Dans le réseau dimensionnel de la mécanique dimensionnelle,
\[
 \mathbf G=-\mathbf S+5\mathbf C+2\mathbf T,
 \qquad [G]=L^3M^{-1}T^{-2}.
\]
Soient
\[
 c_{\rm int}\in\mathcal L_C^+,
 \qquad t_0\in\mathcal L_T^+,
 \qquad \ell_0=c_{\rm int}t_0,
 \qquad \hbar_{\rm int}\in\mathcal L_S^+.
\]

\HMTClaim{MD-R-G-TRANSDUCER}
\begin{definition}[Transducteur gravitationnel]
\label{def:transductor-gravitatorio-rev11}
Pour \(\theta_G>0\), on définit
\begin{equation}
 \boxed{
 G_{\rm int}(\theta_G)
 =\theta_G\frac{c_{\rm int}^{5}t_0^2}{\hbar_{\rm int}}
 =\theta_G\frac{c_{\rm int}^{3}\ell_0^2}{\hbar_{\rm int}}
 \in\mathcal L_G^+.}
 \label{eq:transductor-gravitatorio-rev11}
\end{equation}
\end{definition}

L'égalité des deux expressions et le type de
\eqref{eq:transductor-gravitatorio-rev11} sont exacts. Si
\[
 L_G^2:=\frac{\hbar_{\rm int}G_{\rm int}}{c_{\rm int}^3},
\]
alors
\begin{equation}
 \theta_G=\left(\frac{L_G}{\ell_0}\right)^2.
 \label{eq:thetaG-razon-longitudes-rev11}
\end{equation}
La coordonnée gravitationnelle est donc un rapport quadratique de longueurs
ou une réponse d'holonomie normalisée. La décade
\(10^{-34}\mathcal U_S\) appartient déjà à \(\hbar_{\rm int}\) : elle fixe la section
absolue d'action et ne constitue pas un facteur relatif ajouté à \(G\).

\subsection{Torsion observable et caractères positifs}

Le retour observable de \(27\) pas
\[
 \mathcal K_{27}:=F_{\rm obs}^{27}
\]
satisfait
\[
 \mathcal K_{27}^{4}=F_{\rm obs}^{108}=I.
\]
Dans la coordonnée \(\theta\in\mathbb Z/108\mathbb Z\), l'ordre est
exactement quatre, car
\(\mathcal K_{27}(\theta)=\theta+27\),
\(\mathcal K_{27}^{2}(\theta)=\theta+54\neq\theta\) et
\(\mathcal K_{27}^{4}(\theta)=\theta\).

\HMTClaim{MD-R-G-POSITIVE-CHARACTER-NOGO}
\begin{theorem}[Trivialité du caractère positif observable]
\label{thm:caracter-positivo-G-rev11}
Tout homomorphisme multiplicatif
\[
 \chi:\langle\mathcal K_{27}\rangle
 \longrightarrow\mathbb R_{>0}^{\times}
\]
cumple \(\chi(\mathcal K_{27})=1\).
\end{theorem}

\begin{proof}
La multiplicativité implique
\(\chi(\mathcal K_{27})^4=\chi(I)=1\).
Le seul réel positif dont la puissance quatrième est un est \(1\).
\end{proof}

Le même cycle transporte bien les caractères unitaires
\[
 \chi_j(\mathcal K_{27})=\exp(2\pi ij/4),
 \qquad j=0,1,2,3.
\]
Le quotient observable admet une phase, tandis que sa torsion élimine une
échelle positive multiplicative non triviale. Une continuation multiplicative doit donc
agir sur un relèvement \(\widetilde{\mathcal K}_{27}\) dont la classe dans
l'abélianisation du groupe d'isotropie soit non de torsion, et doit sélectionner
naturellement tant le caractère que sa normalisation.

\subsection{Mémoire chronologique et réponse spectrale}
\label{sec:respuesta-espectral-G-rev11}

Le revêtement chronologique dirigé réalise le retour de \(27\) pas comme la
translation
\[
 (G_{\rm or},S)\longmapsto(G_{\rm or},S)+(1,18),
\]
d'ordre de semi-groupe infini. Cette voie conserve l'accumulation, mais sa
réponse spectrale requiert une interférence entre histoires.

Soit \(\Gamma=(V,E)\) un multigraphe orienté fini muni d'une mesure positive,
de poids \(a_\ast(e)\geq0\) et d'un cocycle \(c(e)\in\mathbb Z^2\). L'opérateur
tordu
\[
 (\mathcal L_{\beta,\vartheta}f)(y)
 =\sum_{e:x\to y}
 e^{-\beta a_\ast(e)}e^{i\langle\vartheta,c(e)\rangle}f(x)
\]
induit l'opérateur positif
\(\mathcal H_{\beta,\vartheta}
=\mathcal L_{\beta,\vartheta}^{*}\mathcal L_{\beta,\vartheta}\).

\HMTClaim{MD-R-G-PHASE-INTERFERENCE}
\begin{proposition}[Critère d'interférence pour la réponse de phase]
\label{prop:interferencia-selector-G-rev11}
Dans une orbite déterministe avec une unique arête entrante par sommet,
\(\mathcal H_{\beta,\vartheta}\) est indépendant de \(\vartheta\). Toute
réponse spectrale de phase non nulle requiert des histoires cohérentes
multiples ou une différentielle magnétique dont l'holonomie subsiste dans les cycles
enrichis.
\end{proposition}

\begin{proof}
Dans l'orbite déterministe, l'unique phase de chaque terme est multipliée par sa
conjuguée dans
\(\mathcal L_{\beta,\vartheta}^{*}\mathcal L_{\beta,\vartheta}\)
et s'annule. L'absence de termes croisés élimine toute dépendance à
\(\vartheta\).
\end{proof}

Si une bande simple \(\lambda_0(\vartheta)\) possède un minimum à l'origine,
sa hessienne \(Q\) est semi-définie positive. Avec
\[
 v_{27}=(1,18),
 \qquad
 \mathcal R_{27}=v_{27}^{\mathsf T}Qv_{27},
\]
la condition \(\mathcal R_{27}\neq0\) constitue le contrôle minimal de
sensibilité à la mémoire accumulée. Une normalisation naturelle
\(\mathfrak N\) permettrait de définir
\[
 \Lambda_G=\ell_0\mathfrak N(Q,v_{27}),
 \qquad
 G_{\rm spec}
 =\frac{c_{\rm int}^{3}}{\hbar_{\rm int}}
  \bigl(\delta_\theta\Lambda_G\bigr)^2,
 \qquad
 \delta_\theta=\frac{A_{\rm deg}-C_{\rm deg}^{\ast}}{360}.
\]
L'homogénéité de ce gabarit est exacte. Sa réalisation physique est
caractérisée par quatre conditions positives : multigraphe enrichi,
réponse non nulle, stabilité par raffinement et indépendance vis-à-vis du changement d'échelle,
avec une normalisation \(\mathfrak N\) fixée avant la confrontation
métrologique.

\subsection{Fonctionnelle géométrique sur l'holonomie enrichie}

La seconde continuation définit directement une fonctionnelle de longueur
\[
 \mathscr L:
 \operatorname{Loop}(\mathcal X_{\rm TPK}^{\rm enr})
 \longrightarrow\mathcal L_L^+\cup\{0\},
\]
invariante par changement cyclique du point de base, conjugaison, réétiquetage
admissible et subdivision. Pour un relèvement enrichi, on écrit
\[
 \Lambda_{27}:=\mathscr L(\widetilde{\mathcal K}_{27}),
 \qquad
 \rho_{27}:=\frac{\Lambda_{27}}{c_{\rm int}t_0},
 \qquad
 \theta_G^{\rm hol}:=(\delta_\theta\rho_{27})^2.
\]
La transduction produit le gabarit homogène
\begin{equation}
 \boxed{
 G_{\rm hol}
 =\frac{c_{\rm int}^{3}}{\hbar_{\rm int}}
  \bigl(\delta_\theta\Lambda_{27}\bigr)^2.}
 \label{eq:plantilla-holonomia-G-rev11}
\end{equation}
La loi de concaténation de cette fonctionnelle sur le cycle de torsion fait
partie de sa définition et de sa preuve. La sélection physique est complétée
en construisant \(\widetilde{\mathcal K}_{27}\) et \(\mathscr L\), en fixant une
normalisation unique et en vérifiant la stabilité de la sortie. La valeur
\(G_{\rm CODATA\,2022}\) intervient exclusivement ensuite, dans la confrontation
métrologique.

\begin{proposition}[Séparation des deux sélecteurs gravitationnels]
\label{prop:separacion-selectores-G}
Le sélecteur circulaire
\(\theta_{108}^{\rm circ}=(54/\pi)^2\), construit par la fermeture
action--Compton--gravité, et le sélecteur holonomique
\(\theta_G^{\rm hol}=(\delta_\theta\rho_{27})^2\) sont deux réalisations
typées distinctes du transducteur \(G_{\rm int}\). Ils ne s'identifient pas l'un à
l'autre et la carte décimale \(\mathscr R_G\) ne constitue pas une évaluation de
\(G_{\rm hol}\).
\end{proposition}
\begin{proof}
Le premier sélecteur dépend uniquement de la fermeture circulaire d'ordre 108. Le
second contient, en outre, la longueur spectrale
\(\Lambda_{27}=\mathscr L(\widetilde{\mathcal K}_{27})\). Une égalité
entre les deux exigerait de calculer \(\Lambda_{27}\) et de prouver la coïncidence des
normalisations ; aucune opération du lecteur décimal n'effectue ce calcul.
Par conséquent, leurs domaines et données d'entrée les distinguent avant toute
confrontation numérique.
\end{proof}

\subsection{Critère de sélection structurelle}

\HMTClaim{MD-R-G-SELECTOR-CRITERION}
\begin{criterion}[Sélecteur structurel de la constante gravitationnelle]
\label{crit:selector-G-rev11}
Un sélecteur HMT de \(G\) est identifié lorsqu'il satisfait conjointement :
\begin{enumerate}
 \item domaine enrichi, composition et orientation déclarés ;
 \item descente à une classe abélienne non de torsion, ou fonctionnelle géométrique
 invariante par point de base, conjugaison, réétiquetage et subdivision ;
 \item réponse positive non triviale et stable sous raffinement ;
 \item normalisation unique et indépendante de l'échelle globale de l'opérateur ;
 \item sélection interne de la coordonnée et de sa position décimale ;
 \item gel du code et des entrées avant d'ouvrir la valeur de
 confrontation ;
 \item réalisation finale comme section de \(\mathcal L_G\) dans une carte
 dimensionnelle commune.
\end{enumerate}
\end{criterion}

Les rapports adimensionnels
\[
 \frac{\mathcal G_{\rm tor}}{\mathcal I_{\rm tor}}=1,
 \qquad
 \frac{\mathcal G_{\rm av}}{\mathcal I_{\rm av}}
 =\frac{\pi_{\rm HMT}}{\varphi^2}
\]
décrivent des lecteurs gravitationnels--inertiels distincts et demeurent
compatibles avec une seule section dimensionnelle \(G\). La construction de cette
section utilise en outre
\((\hbar_{\rm int},c_{\rm int},\ell_0)\) et le sélecteur \(\theta_G\).

\subsection{Évaluation de la carte décimale gravitationnelle et confrontation métrologique ultérieure}

La représentation décimale conservée dans le corpus est
\begin{equation}
 \mathscr R_G([n,k,d]_3;t,e)
 :=10^{-n}\left(k+\frac{t}{10^3}+\frac{e}{10^d}\right).
 \label{eq:lector-decimal-G-rev11}
\end{equation}
Pour la carte
\([n,k,d]_3=[11,6,5]_3\), \(t=674\) et \(e=30\), l'évaluation rationnelle est
\begin{equation}
 \boxed{
 \mathscr R_G
 =10^{-11}\left(6+\frac{674}{1000}+\frac{30}{10^5}\right)
 =6.67430\times10^{-11}.}
 \label{eq:evaluacion-decimal-G-rev11}
\end{equation}
Le triplet, les deux blocs et la position décimale constituent les données
déclarées de cette carte. Son évaluation est exacte ; sa sélection structurelle
est régie par le \cref{crit:selector-G-rev11}.

Après choix de la base SI de \(\mathcal L_G\), la confrontation en vigueur est
\begin{equation}
 \boxed{
 G_{\rm CODATA\,2022}=6.67430(15)\times10^{-11}
 \ \mathrm{m^3\,kg^{-1}\,s^{-2}}.}
 \label{eq:G-CODATA-2022-rev11}
\end{equation}

\begin{center}
\begin{tabularx}{0.94\textwidth}{@{}L{0.27\textwidth}Y@{}}
\toprule
Champ & Déclaration métrologique\\
\midrule
Valeur centrale &
\(6.67430\times10^{-11}\ \mathrm{m^3\,kg^{-1}\,s^{-2}}\)\\
Incertitude type &
\(u(G)=0.00015\times10^{-11}
=1.5\times10^{-15}\ \mathrm{m^3\,kg^{-1}\,s^{-2}}\)\\
Incertidumbre relativa & \(u_r(G)\simeq2.2\times10^{-5}\)\\
Source et temporalité &
Valeurs recommandées CODATA 2022, publiées en 2024
\cite{CODATA2022}\\
Fonction probatoire &
Confrontation externe ultérieure ; l'unité et l'incertitude n'interviennent pas dans
l'évaluation rationnelle de \eqref{eq:evaluacion-decimal-G-rev11}.\\
\bottomrule
\end{tabularx}
\end{center}

\paragraph{Reproductibilité du lecteur gravitationnel.}
Une énumération indépendante vérifie conjointement le type dimensionnel, le
lecteur décimal et les identités de monodromie sur le domaine fini
déclaré.

\begin{proposition}[Chaîne causale de la constante gravitationnelle]
L'action covariante, le hamiltonien contraint et l'équation
Einstein--Schrödinger gouvernent la dynamique gravitationnelle. Le type dimensionnel
et le transducteur \eqref{eq:transductor-gravitatorio-rev11} sont exacts ;
l'évaluation \eqref{eq:evaluacion-decimal-G-rev11} est exacte pour sa carte
déclarée. La valeur, l'unité et l'incertitude de
\eqref{eq:G-CODATA-2022-rev11} constituent une confrontation externe CODATA 2022. La
sélection autonome de \(\theta_G\) est typée par le relèvement
enrichi, la réponse interférente et la fonctionnelle d'holonomie soumis
au \cref{crit:selector-G-rev11}.
\end{proposition}
